% ------------------------------------------------------------------------
% file `3_uaa1_fonctions_2022-2023-transferer-6-solution-body.tex'
%   in folder `xsim/'
%
%     solution of type `transferer' with id `6'
%
% generated by the `soltransferer' environment of the
%   `xsim' package v0.21 (2022/02/12)
% from source `3_uaa1_fonctions_2022-2023' on 2022/11/05 on line 292
% ------------------------------------------------------------------------
    1 point pour la réponse, 1 point pour le raisonnement.

    Déterminons les expressions analytiques des deux fonctions.

    \textbf{\faVial~Éprouvette A}

    Les points $(0,12)$ et $(2,11)$ appartiennent à $f_A$ par lecture graphique.\par
    \begin{tblr}[t]{columns={leftsep=0pt},
            row{1}={font=\bfseries},
            colspec={XX},
            cell{3}{1}={r=1,c=2}{c}
        }
        Déterminons m                     & Déterminons p                    \\
        $m=-\dfrac{1}{2}$                 & $(0,12)\in f_A \Rightarrow p=12$ \\
        \fbox{$f_A(x)=-\dfrac{1}{2}x+12$} &
    \end{tblr}

    \textbf{\faVial~Éprouvette B}

    Les points $(0,8)$ et $(4,7)$ appartiennent à $f_B$ par lecture graphique.\par
    \begin{tblr}[t]{columns={leftsep=0pt},
            row{1}={font=\bfseries},
            colspec={XX},
            cell{3}{1}={r=1,c=2}{c}
        }
        Déterminons m                    & Déterminons p                  \\
        $m=-\dfrac{1}{4}$                & $(0,8)\in f_A \Rightarrow p=8$ \\
        \fbox{$f_B(x)=-\dfrac{1}{4}x+8$} &
    \end{tblr}

    \begin{enumerate}
        \item Trouvons la racine de $f_A$,
              \begin{align*}
                  0   & =-\dfrac{1}{2}x+12 \\
                  -12 & = -\dfrac{1}{2}x   \\
                  x   & = 24
              \end{align*}
              L'évaporation de l'huile contenu dans l'éprouvette A sera complète après 24 jours.
        \item Trouvons l'abscisse du point d'intersection des deux fonctions,
              \begin{align*}
                  -\dfrac{1}{2}x+12              & = -\dfrac{1}{4}x+8 \\
                  -\dfrac{1}{2}x + \dfrac{1}{4}x & = 8 -12            \\
                  -\dfrac{1}{4}x                 & = -4               \\
                  x                              & = 16
              \end{align*}
              L'huile contenue dans les deux éprouvettes sera identique après 16 jours.
    \end{enumerate}
